\documentclass[11pt]{article}
\usepackage[margin=25mm]{geometry}
\usepackage{fontspec}
\setmainfont{Latin Modern Roman}
\usepackage{amsmath,amssymb,mathtools}
\usepackage{unicode-math}
\setmathfont{Latin Modern Math}
\usepackage{xurl,hyperref}
\hypersetup{colorlinks=true,urlcolor=blue}
\setcounter{secnumdepth}{0}
\setlength{\parindent}{0pt}
\setlength{\parskip}{5pt}
\setlength{\emergencystretch}{3em}
\newcommand{\sref}[2]{\hyperlink{source:#1:#2}{[S#2]}}
\newcommand{\cataloguescope}[1]{\paragraph{Recorded target.}#1}
\newcommand{\R}{\mathbb R}
\newcommand{\C}{\mathbb C}
\newcommand{\N}{\mathbb N}
\newcommand{\Z}{\mathbb Z}
\newcommand{\Q}{\mathbb Q}
\newcommand{\F}{\mathbb F}
\newcommand{\eps}{\varepsilon}
\begin{document}
\section{Hauptmodul differential equations for an even-weight mixed-mock class}\label{statement}

\subsection{Definitions and mathematical statement}

\paragraph{Shared notation.}$\mathbb N=\{1,2,\ldots\}$; the standard number fields have their usual meanings. All additional conventions are stated locally.


\paragraph{Definitions and notation.}A modular function is an invariant meromorphic function for a specified modular group. Differentiation with respect to a local modular coordinate $t(\tau)$ is $(dt/d\tau)^{-1}d/d\tau$ away from its critical points. A differential equation may be linear or nonlinear and can have different coefficient classes. Mixed mock forms have an additional nonholomorphic-completion structure.

\paragraph{Catalogue wording retained for traceability.}
Do mixed mock modular forms satisfy differential equations with respect to a modular function?
\paragraph{Scope and source qualification.}The source fixes neither a modular group nor an order or coefficient field for the equation. Choosing an arbitrary equation constructed from a given function does not establish the structural property sought. The catalogue title/status was paired with a machine-generated research attempt; that attempt is not peer-reviewed evidence of settlement. The statement covers the original source questions rather than only the title’s subcase.
\subsection{Short English statement}
Find differential equations for mixed mock modular forms when a modular function is used as the variable.
\subsection{Research attempt}
\paragraph{Model, date and investigated scope.}
OpenAI GPT-6 Astra (configured model \texttt{gpt-6-astra}), ultra reasoning; generated 11 October 2026. Derived the exact coordinate-change formula for a second-order ODE and a branch obstruction to coefficient descent. No structural mixed-mock equation over a modular function field is established.

\paragraph{Sources and literature check.}
Read the supplied complete definition, catalogue wording, scope qualifications and source list for this canonical ID. The original-source citations are inherited as supplied references; no fresh primary-source or novelty search was performed for this attempt. No current open-status claim is inferred from this calculation.

\paragraph{Changing the independent variable rigorously.}Let \(t=t(\tau)\) be holomorphic with \(t'(\tau)\ne0\) on a simply connected local chart, and suppose \(f\) satisfies a known equation \(f''+a(\tau)f'+b(\tau)f=0\). Write \(f(\tau)=y(t(\tau))\). Chain-rule differentiation gives
\[y_{tt}+\left(\frac{t''}{(t')^2}+\frac a{t'}\right)y_t+\frac b{(t')^2}y=0.
\]
This is a differential equation in the modular coordinate only if the two displayed coefficients are single-valued functions of \(t\) in the intended global setting. A local inverse alone gives no such arithmetic or meromorphic coefficient property.
\paragraph{A branch countercheck.}Take \(t=\tau^2\) away from zero and the equation \(f''+f'=0\). Its transformed first-derivative coefficient is \(1/(2t)+1/(2\tau)\). The final term changes sign between the two lifts \(\tau\) and \(-\tau\). Thus a perfectly regular differential equation upstairs need not descend to a scalar equation with coefficients on the \(t\)-plane. Critical points where \(t'=0\) create further singularities and require separate local analysis.
\paragraph{Gap.}This example is not claimed to be a mock modular form; it isolates the actual descent obstruction in a tempting chain-rule argument. For mixed mock forms one must first derive a non-tautological differential system from their modular completion/shadow, then prove coefficient descent in a specified modular function field, including branch points and cusps. No such system is obtained here. Merely writing \(dy/dt\) using a local inverse does not answer the source's structural question.

\paragraph{Verification and adversarial audit.}
Checked both chain-rule coefficients, two-sheet sign counterexample, critical-point exclusion and non-mock status of the diagnostic example. This was an LLM adversarial audit, not independent human verification; no theorem prover was used. Prompts A then B applied to this ID: Checked both chain-rule coefficients, two-sheet sign counterexample, critical-point exclusion and non-mock status of the diagnostic example. Scope remains limited: Derive a completion-based differential system and prove its coefficients descend globally to the chosen modular function field.

Coordinator checked the change-of-variable coefficients and sign-changing two-sheet example obstructing global descent. This additional check was model-assisted, not independent human verification.

\paragraph{Tools, process and contributions.}
Python 3 and local shell used to read supplied files and serialize this manuscript; mathematics derived symbolically in the text, with no numerical experiment or formal proof assistant. Contribution prompt A was followed by prompt B. GPT-6 Astra with ultra reasoning drafted the research and performed its stated self-audit; any additional coordinator contribution is identified in the verification and audit. A GPT-6-based Codex coordinator assembled the manuscript and checked repository integration; its exact variant and reasoning configuration were not exposed. No human mathematical contribution, independent human proof check or formal proof certification is supplied. Complete derivations and the consulted source references are included here; private reasoning transcripts are not published.

\paragraph{Final status and first remaining gap.}
\textbf{UNRESOLVED}. Derive a completion-based differential system and prove its coefficients descend globally to the chosen modular function field.

\section{Completion Estimate}
\noindent\textbf{COMPLETION: 0\%}
The calculation is an elementary benchmark or conditional reduction; it establishes no claimed new portion of the complete source target.

\subsection{Sources}
\begin{itemize}
\item[\hypertarget{source:statement:1}{S1}] ULAM.AI and the UnsolvedMath Contributors, UnsolvedMath, catalogue entry. Catalogue curation under CC BY 4.0; definitions and English exposition added locally.
\url{https://huggingface.co/datasets/ulamai/UnsolvedMath}.
\item[\hypertarget{source:statement:2}{S2}] American Institute of Mathematics, Gromov-Witten invariants and number theory, source item 5.2.
\url{http://aimpl.org/gromwitnumthry/5/}.
\end{itemize}
\label{end:statement}
\end{document}
